\documentclass[11pt]{article}

\usepackage[utf8]{inputenc}
\usepackage[T1]{fontenc}
\usepackage{lmodern}
\usepackage{geometry}
\usepackage{amsmath,amssymb,amsfonts,amsthm,mathtools}
\usepackage{bm}
\usepackage{enumitem}
\usepackage{microtype}
\usepackage{hyperref}
\usepackage{titlesec}
\usepackage{booktabs}
\usepackage{array}
\usepackage{setspace}
\usepackage{mathrsfs}
\usepackage{physics}

\geometry{margin=1in}
\hypersetup{colorlinks=true, linkcolor=black, urlcolor=blue, citecolor=black}
\setlist[enumerate]{topsep=0.4em,itemsep=0.35em}
\setlist[itemize]{topsep=0.3em,itemsep=0.25em}
\titleformat{\section}{\large\bfseries}{\thesection.}{0.5em}{}
\titleformat{\subsection}{\normalsize\bfseries}{\thesubsection.}{0.5em}{}
\setstretch{1.06}

\newcommand{\Aether}{\AE{}ther}
\newcommand{\Aflow}{\AE{}ther-flow}
\newcommand{\TheoryName}{\Aflow{} Interpretation of Relativity}
\newcommand{\TheoryProgram}{\Aether{} / \Aflow{} framework}

\newcommand{\CompletedMetric}{g^{\sharp}}
\newcommand{\MatterSpace}{\Psi}

\newcommand{\Car}{\mathrm{car}}
\newcommand{\Cur}{\mathrm{cur}}
\newcommand{\Hom}{\mathrm{hom}}

\newcommand{\ForSpace}[1]{\mathcal I_{#1}^{\mathrm{for}}}
\newcommand{\PrimShell}{\mathcal S_{\mathrm{prim}}}
\newcommand{\MetricOut}{\mathscr G_{\mathrm{out}}}
\newcommand{\MatterOut}{\mathscr T_{\mathrm{out}}}

\newcommand{\ProtoSectionPackage}{\mathfrak Q^{\mathrm{sub,proto,secgen}}}
\newcommand{\ProtoSectionBody}[1]{\mathcal B_{#1}^{\mathrm{psec}}}
\newcommand{\ProtoSectionShell}[1]{\mathcal Z_{#1}^{\mathrm{psec}}}

\newcommand{\PreProtoSectionPackage}{\mathfrak R^{\mathrm{sub,proto,presec}}}
\newcommand{\PreProtoSectionMinSectionCore}{\mathfrak R^{\mathrm{sub,proto,presec,sec}}}
\newcommand{\PreProtoSectionResidualPackage}{\mathfrak R^{\mathrm{sub,proto,presec,res}}}

\newcommand{\PrePreProtoSectionPullbackCore}{\mathfrak S^{\mathrm{sub,proto,pppresec,pb}}}
\newcommand{\PrePreProtoSectionVertical}[1]{V_{#1}^{\mathrm{pppresec}}}
\newcommand{\PrePreProtoSectionResidualBody}[1]{\mathcal D_{#1}^{\mathrm{pppresec}}}
\newcommand{\PrePreProtoSectionBodyResponse}[1]{\mathscr R_{#1}^{\mathrm{pppresec,pb,body}}}

\newcommand{\PullbackOperatorCore}{\mathfrak S^{\mathrm{sub,proto,pppresec,pb,op}}}
\newcommand{\PullbackBodyOperatorCore}{\mathfrak S^{\mathrm{sub,proto,pppresec,pb,bodyop}}}
\newcommand{\PullbackBodyResponseCore}{\mathfrak S^{\mathrm{sub,proto,pppresec,pb,bodyresp}}}

\newcommand{\ProjectionOperatorCore}{\mathfrak S^{\mathrm{disp,resp,projop}}}
\newcommand{\CurrentBodyImageCore}{\mathfrak S^{\mathrm{disp,resp,cbimg}}}
\newcommand{\SubstrateCarrier}[1]{U_{#1}^{\mathrm{disp,resp}}}
\newcommand{\SubstrateBody}[1]{\mathcal U_{#1}^{\mathrm{disp,resp}}}
\newcommand{\SubstrateProjection}[1]{P_{#1}^{\mathrm{disp,resp}}}
\newcommand{\SubstrateOperator}[1]{K_{#1}^{\mathrm{disp,resp}}}
\newcommand{\SubstrateKernel}[1]{N_{#1}^{\mathrm{disp,resp}}}
\newcommand{\SubstrateEnergy}[1]{\mathscr E_{#1}^{\mathrm{disp,resp}}}
\newcommand{\CanonicalLift}[1]{J_{#1}^{\mathrm{disp,resp,can}}}
\newcommand{\SchurResponse}[1]{M_{#1}^{\mathrm{disp,resp}}}
\newcommand{\CoreForcedBodyResponse}[1]{\widehat{\mathscr R}_{#1}^{\mathrm{disp,resp,projop}}}
\newcommand{\CanonicalImageBody}[1]{\mathcal U_{#1}^{\mathrm{disp,resp,cbimg}}}
\newcommand{\CanonicalImageShell}[1]{\mathcal Z_{#1}^{\mathrm{disp,resp,cbimg}}}
\newcommand{\CanonicalImageResponse}[1]{\widehat{\mathscr R}_{#1}^{\mathrm{disp,resp,cbimg}}}

\newtheorem{definition}{Definition}
\newtheorem{proposition}{Proposition}
\newtheorem{remark}{Remark}
\newtheorem{corollary}{Corollary}
\newtheorem{theorem}{Theorem}

\title{\textbf{The \TheoryName{}}\\[0.4em]
\large Primitive-Reservoir Observer-Reduction\\
\large Section-Centered Hidden-Response\\
\large Projection-Operator Core\\
\large Collapse to a Canonical Current-Body Image Core}
\author{}
\date{}

\begin{document}

\maketitle

\begin{abstract}
The current active-primary theorem already proves that the full section-centered displacement-response substrate law carries no benchmark-facing freedom beyond a smaller current-body-realizing projection-operator core \(\ProjectionOperatorCore\): once the substrate body, the surjective displacement projection, the coercive positive self-adjoint substrate-response operator, the exact-shell discipline, the current-body realization, and benchmark faithfulness are fixed, the \(K\)-orthogonal minimal lift is uniquely forced. The live burden therefore no longer sits on the full substrate law or on the separately primitive lift family.

The present note asks the next honest internal question strictly inside that smaller core. Does the full substrate body itself still carry independent benchmark-facing freedom, or do all current uses already factor through the canonical image of the fixed current displacement body under the forced minimal lift? The answer proved here is that the larger body is still presentationally redundant at benchmark scope. Once the current displacement body, the surjective projection, the coercive response operator, and benchmark faithfulness are fixed, every point of the full substrate body decomposes uniquely into the canonical minimal-lift representative of the same current displacement together with a kernel direction, and that kernel direction contributes only nonnegative excess response energy without altering the forced body-response core or any downstream benchmark-facing consequence. What remains live is the smaller canonical current-body image core \(\CurrentBodyImageCore\).

The benchmark boundary remains fixed throughout. The one operative metric is still exactly \(\CompletedMetric\), the ordinary matter route is unchanged, there is no second operative metric, the low-energy matter law is not enlarged, and no stronger promotion language is licensed. The positive result is therefore a genuine smaller theorem strictly inside the current projection-operator-core frontier.
\end{abstract}

\section{Introduction}

The active derivational program of \TheoryName{} has again reached the point where another deeper relay would no longer be the honest next step. The current section-centered hidden-response theorem already compressed the full displacement-response substrate law to the smaller current-body-realizing projection-operator core \(\ProjectionOperatorCore\) by proving that the \(K\)-orthogonal minimal-lift family is forced. After that gain, another manuscript that merely restates the same projection core through a deeper presentation would not remove the remaining benchmark-facing freedom.

The next honest question is therefore internal to that smaller core itself. Relative to the already fixed current displacement body and the already forced canonical lift, is the entire full substrate body still independently live, or do current benchmark-facing uses see only the canonical image of the current body inside the carrier? The present note proves the latter. The benchmark-relevant representative above each current displacement is the canonical minimal-lift point. Body directions beyond that image are purely kernel directions, and at current scope they do not change the forced body-response core or any downstream benchmark-facing consequence.

\paragraph{Framework and claim boundary.}
\TheoryProgram{} denotes the broader \Aether{} / \Aflow{} framework built on the claims that the \Aether{} is the underlying four-dimensional substrate of reality, the \Aflow{} is its intrinsic ordered motion, observed three-dimensional space is the local experiential slice of that deeper substrate, S-time is the experienced order of change arising from matter, light, and the \Aflow{}, and observed expansion is the three-dimensional appearance of deeper four-dimensional ordered motion. Within that framework, \TheoryName{} denotes the exact-closure theory in which the effective gravitational dynamics are adopted to be exactly Einsteinian with universal matter coupling. In the active sequence, \emph{adoption} means use of that established relativistic dynamics without claiming substrate derivation, while \emph{derivation} is reserved for a first-principles recovery from explicit substrate variables.

The present note stays strictly inside that boundary. It does not introduce a second operative metric, it does not enlarge the low-energy matter law, and it does not reopen the full displacement-response substrate law or the larger pullback-body-operator, pullback-operator, pullback-quadratic, hidden-response, residual-normal-form, or pre-proto-section presentations as if they were again independently live. Its narrower benchmark-facing role is to identify whether the full substrate-body presentation still matters once the current displacement body and the forced canonical minimal lift are already fixed.

\section{Fixed Inputs}

\begin{definition}[Recorded primitive continuation data]
\label{def:fixed}
Fix the following continuation data from the active primitive-reservoir observer-reduction line.
\begin{enumerate}
    \item For each sector label \(\sigma\in\{\Car,\Cur,\Hom\}\), a primitive forcing shell
    \[
    \ForSpace{\sigma}.
    \]
    \item The product shell
    \[
    \PrimShell
    \equiv
    \ForSpace{\Car}\times \ForSpace{\Cur}\times \ForSpace{\Hom}.
    \]
    \item The recorded metric and ordinary-matter outputs
    \[
    \MetricOut:\PrimShell\to\{\text{observer metrics}\},
    \qquad
    \MatterOut:\PrimShell\times\MatterSpace\to\{\text{observer stress tensors}\},
    \]
    satisfying
    \begin{equation}
    \MetricOut(\mathbf w)=\CompletedMetric,
    \qquad
    \MatterOut(\mathbf w,\psi)
    =
    T_{\mu\nu}^{\mathrm{matter}}[\CompletedMetric,\psi]
    \label{eq:fixedoutputs}
    \end{equation}
    for every \(\mathbf w\in\PrimShell\) and every matter configuration \(\psi\in\MatterSpace\).
\end{enumerate}
\end{definition}

\begin{definition}[Recorded proto-section benchmark base and current body-response datum]
\label{def:bodyresp}
Fix \(\sigma\in\{\Car,\Cur,\Hom\}\). The active line already fixes:
\begin{enumerate}
    \item the current proto-section benchmark base \(\ProtoSectionPackage\), in particular the proto-section body \(\ProtoSectionBody{\sigma}\) and shell \(\ProtoSectionShell{\sigma}\);
    \item the current section-centered displacement body
    \[
    \PrePreProtoSectionResidualBody{\sigma}
    \subset
    \ProtoSectionBody{\sigma}\times\PrePreProtoSectionVertical{\sigma}
    \]
    whose fiber
    \[
    \PrePreProtoSectionResidualBody{\sigma}(y)
    \equiv
    \left\{
    v\in\PrePreProtoSectionVertical{\sigma}:(y,v)\in\PrePreProtoSectionResidualBody{\sigma}
    \right\}
    \]
    is nonempty, compact, convex, contains \(0\), and has nonempty interior for every \(y\in\ProtoSectionBody{\sigma}\);
    \item the body-restricted response
    \[
    \PrePreProtoSectionBodyResponse{\sigma}:
    \PrePreProtoSectionResidualBody{\sigma}\to[0,\infty),
    \]
    smooth in \(y\), fiberwise quadratic in \(v\), and satisfying
    \[
    \PrePreProtoSectionBodyResponse{\sigma}(y,v)=0
    \qquad\Longleftrightarrow\qquad
    v=0;
    \]
    \item benchmark faithfulness, meaning preservation of \eqref{eq:fixedoutputs}, no second operative metric, no enlarged low-energy matter law, no change to the ordinary matter route, and no premature promotion from adoption to completed derivation.
\end{enumerate}
\end{definition}

\begin{definition}[Recorded downstream chain below the current projection frontier]
\label{def:downstream}
Fix \(\sigma\in\{\Car,\Cur,\Hom\}\). The active line already records that, once a benchmark-faithful current-body-realizing projection-operator core over the fixed proto-section benchmark base is available, the current observer-relevant pullback body-response core \(\PullbackBodyResponseCore\), the current pullback body-operator core \(\PullbackBodyOperatorCore\), the current pullback-operator core \(\PullbackOperatorCore\), the current pullback-quadratic core \(\PrePreProtoSectionPullbackCore\), the current section-centered residual-normal-form realization \(\PreProtoSectionResidualPackage\), the current benchmark-faithful pre-proto-section package \(\PreProtoSectionPackage\), the current minimizer-section core \(\PreProtoSectionMinSectionCore\), and the previously recorded continuity-Hessian compressed core and benchmark one-metric observer package follow unchanged.
\end{definition}

\section{Current Projection-Operator-Core Frontier}

\begin{definition}[Current-body-realizing section-centered displacement-response projection-operator core]
\label{def:core}
Fix \(\sigma\in\{\Car,\Cur,\Hom\}\). A \emph{benchmark-faithful current-body-realizing section-centered displacement-response projection-operator core} \(\ProjectionOperatorCore\) for sector \(\sigma\) consists of:
\begin{enumerate}
    \item a finite-dimensional Euclidean substrate carrier
    \[
    \SubstrateCarrier{\sigma};
    \]
    \item a compact convex section-centered substrate body
    \[
    \SubstrateBody{\sigma}
    \subset
    \ProtoSectionBody{\sigma}\times\SubstrateCarrier{\sigma}
    \]
    whose fiber
    \[
    \SubstrateBody{\sigma}(y)
    \equiv
    \left\{
    u\in\SubstrateCarrier{\sigma}:(y,u)\in\SubstrateBody{\sigma}
    \right\}
    \]
    is nonempty, compact, convex, contains \(0\), and has nonempty interior for every \(y\in\ProtoSectionBody{\sigma}\);
    \item a smooth family of surjective linear displacement projections
    \[
    \SubstrateProjection{\sigma}(y):
    \SubstrateCarrier{\sigma}\to\PrePreProtoSectionVertical{\sigma},
    \qquad
    y\in\ProtoSectionBody{\sigma};
    \]
    \item a smooth family of positive self-adjoint substrate-response operators
    \[
    \SubstrateOperator{\sigma}(y):
    \SubstrateCarrier{\sigma}\to\SubstrateCarrier{\sigma},
    \qquad
    y\in\ProtoSectionBody{\sigma},
    \]
    satisfying the uniform coercive bound
    \begin{equation}
    \ev{u,\SubstrateOperator{\sigma}(y)u}
    \ge
    \mu_{\sigma}\,\|u\|^{2}
    \qquad
    \text{for all }
    y\in\ProtoSectionBody{\sigma},
    \ \ u\in\SubstrateCarrier{\sigma}
    \label{eq:coercive}
    \end{equation}
    for some constant \(\mu_{\sigma}>0\);
    \item exact-shell discipline:
    \[
    \left\{
    y\in\ProtoSectionBody{\sigma}:
    \exists u\in\SubstrateBody{\sigma}(y)
    \text{ with }
    \SubstrateProjection{\sigma}(y)u=0
    \right\}
    =
    \ProtoSectionShell{\sigma};
    \]
    \item current-body realization:
    \begin{equation}
    \PrePreProtoSectionResidualBody{\sigma}(y)
    =
    \SubstrateProjection{\sigma}(y)\bigl(\SubstrateBody{\sigma}(y)\bigr)
    \qquad
    \text{for every }
    y\in\ProtoSectionBody{\sigma};
    \label{eq:realizing}
    \end{equation}
    \item benchmark faithfulness in the sense of Definition \ref{def:bodyresp}(4).
\end{enumerate}
\end{definition}

\begin{remark}
Definition \ref{def:core} is the full current frontier after the minimal-lift collapse theorem has already spent the separately primitive lift family. The present note asks whether the full substrate body still carries independent benchmark-facing freedom once the already-fixed current displacement body and the now-forced canonical lift are held fixed.
\end{remark}

\section{Canonical Lift and Forced Compatibility}

\begin{proposition}[Positive Schur-response operator]
\label{prop:schur}
Fix a current-body-realizing projection-operator core in the sense of Definition \ref{def:core}. For each \(y\in\ProtoSectionBody{\sigma}\), define
\[
\SchurResponse{\sigma}(y)
\equiv
\SubstrateProjection{\sigma}(y)
\SubstrateOperator{\sigma}(y)^{-1}
\SubstrateProjection{\sigma}(y)^*:
\PrePreProtoSectionVertical{\sigma}\to\PrePreProtoSectionVertical{\sigma},
\]
where \((\cdot)^*\) denotes Euclidean adjoint with respect to the fixed inner products on \(\SubstrateCarrier{\sigma}\) and \(\PrePreProtoSectionVertical{\sigma}\). Then \(\SchurResponse{\sigma}(y)\) is positive self-adjoint and invertible for every \(y\in\ProtoSectionBody{\sigma}\).
\end{proposition}

\begin{proof}
Self-adjointness is immediate from the self-adjointness of \(\SubstrateOperator{\sigma}(y)^{-1}\). For any \(0\neq w\in\PrePreProtoSectionVertical{\sigma}\),
\[
\ev{w,\SchurResponse{\sigma}(y)w}
=
\ev{
\SubstrateOperator{\sigma}(y)^{-1/2}\SubstrateProjection{\sigma}(y)^*w,
\SubstrateOperator{\sigma}(y)^{-1/2}\SubstrateProjection{\sigma}(y)^*w
}
\ge 0.
\]
If this quantity vanishes, then \(\SubstrateProjection{\sigma}(y)^*w=0\). Because \(\SubstrateProjection{\sigma}(y)\) is surjective, its Euclidean adjoint is injective, so \(w=0\), a contradiction. Thus \(\SchurResponse{\sigma}(y)\) is positive definite and therefore invertible on the finite-dimensional space \(\PrePreProtoSectionVertical{\sigma}\).
\end{proof}

\begin{definition}[Canonical lift and core-forced response]
\label{def:canonical}
Fix a current-body-realizing projection-operator core. For each \(y\in\ProtoSectionBody{\sigma}\), define the \emph{canonical minimal lift}
\begin{equation}
\CanonicalLift{\sigma}(y)
\equiv
\SubstrateOperator{\sigma}(y)^{-1}\SubstrateProjection{\sigma}(y)^*
\SchurResponse{\sigma}(y)^{-1}
:
\PrePreProtoSectionVertical{\sigma}\to\SubstrateCarrier{\sigma}.
\label{eq:canonical}
\end{equation}
Define the quadratic substrate-response energy by
\[
\SubstrateEnergy{\sigma}(y,u)
\equiv
\frac{1}{2}
\ev{u,\SubstrateOperator{\sigma}(y)u}
\]
and the \emph{core-forced body response}
\begin{equation}
\CoreForcedBodyResponse{\sigma}(y,v)
\equiv
\frac{1}{2}
\ev{
\CanonicalLift{\sigma}(y)v,
\SubstrateOperator{\sigma}(y)\CanonicalLift{\sigma}(y)v
}
\qquad
\text{for }
(y,v)\in\PrePreProtoSectionResidualBody{\sigma}.
\label{eq:coreresponse}
\end{equation}
\end{definition}

\begin{proposition}[Canonical lift properties]
\label{prop:canonical}
Fix a current-body-realizing projection-operator core. Then, for every \(y\in\ProtoSectionBody{\sigma}\),
\begin{enumerate}
    \item
    \[
    \SubstrateProjection{\sigma}(y)\circ\CanonicalLift{\sigma}(y)
    =
    \mathrm{id}_{\PrePreProtoSectionVertical{\sigma}};
    \]
    \item
    \[
    \ev{
    \CanonicalLift{\sigma}(y)v,
    \SubstrateOperator{\sigma}(y)k
    }
    =
    0
    \]
    for every \(v\in\PrePreProtoSectionVertical{\sigma}\) and every \(k\in\SubstrateKernel{\sigma}(y)\equiv\ker\SubstrateProjection{\sigma}(y)\);
    \item every \(u\in\SubstrateCarrier{\sigma}\) admits the unique decomposition
    \[
    u
    =
    \CanonicalLift{\sigma}(y)\bigl(\SubstrateProjection{\sigma}(y)u\bigr)
    +
    k,
    \qquad
    k\in\SubstrateKernel{\sigma}(y);
    \]
    \item \(\CanonicalLift{\sigma}(y)\) is linear in the displacement argument and smooth in \(y\).
\end{enumerate}
\end{proposition}

\begin{proof}
Item (1) follows directly from Definition \ref{def:canonical}:
\[
\SubstrateProjection{\sigma}(y)\CanonicalLift{\sigma}(y)
=
\SchurResponse{\sigma}(y)\SchurResponse{\sigma}(y)^{-1}
=
\mathrm{id}.
\]
For item (2), let \(k\in\SubstrateKernel{\sigma}(y)\). Then
\[
\ev{
\CanonicalLift{\sigma}(y)v,
\SubstrateOperator{\sigma}(y)k
}
=
\ev{
\SubstrateProjection{\sigma}(y)^*
\SchurResponse{\sigma}(y)^{-1}v,
k
}
=
\ev{
\SchurResponse{\sigma}(y)^{-1}v,
\SubstrateProjection{\sigma}(y)k
}
=
0.
\]
For item (3), if \(u\in\SubstrateCarrier{\sigma}\), set
\[
k
\equiv
u-\CanonicalLift{\sigma}(y)\bigl(\SubstrateProjection{\sigma}(y)u\bigr).
\]
Applying \(\SubstrateProjection{\sigma}(y)\) and using item (1) gives \(\SubstrateProjection{\sigma}(y)k=0\), so \(k\in\SubstrateKernel{\sigma}(y)\). Uniqueness follows because item (1) fixes the non-kernel part.

Linearity in the displacement argument is immediate from \eqref{eq:canonical}. Smoothness in \(y\) follows because \(\SubstrateProjection{\sigma}(y)\), \(\SubstrateOperator{\sigma}(y)\), their adjoints and inverses, and \(\SchurResponse{\sigma}(y)^{-1}\) all vary smoothly with \(y\).
\end{proof}

\begin{proposition}[Forced exact-shell/current-body compatibility]
\label{prop:compat}
Under Definition \ref{def:core}, the projected zero slice of the current substrate body is exactly the fixed proto-section shell:
\[
\left\{
y\in\ProtoSectionBody{\sigma}:0\in\PrePreProtoSectionResidualBody{\sigma}(y)
\right\}
=
\ProtoSectionShell{\sigma}.
\]
\end{proposition}

\begin{proof}
If \(0\in\PrePreProtoSectionResidualBody{\sigma}(y)\), then by the current-body realization \eqref{eq:realizing} there exists \(u\in\SubstrateBody{\sigma}(y)\) with \(\SubstrateProjection{\sigma}(y)u=0\). Definition \ref{def:core}(5) then gives \(y\in\ProtoSectionShell{\sigma}\). Conversely, if \(y\in\ProtoSectionShell{\sigma}\), Definition \ref{def:core}(5) yields \(u\in\SubstrateBody{\sigma}(y)\) with \(\SubstrateProjection{\sigma}(y)u=0\). Using \eqref{eq:realizing} again shows \(0\in\PrePreProtoSectionResidualBody{\sigma}(y)\).
\end{proof}

\begin{proposition}[Energy split over the canonical image]
\label{prop:split}
Fix a current-body-realizing projection-operator core. For every \(y\in\ProtoSectionBody{\sigma}\) and every \(u\in\SubstrateCarrier{\sigma}\), if
\[
v\equiv\SubstrateProjection{\sigma}(y)u,
\qquad
k\equiv u-\CanonicalLift{\sigma}(y)v\in\SubstrateKernel{\sigma}(y),
\]
then
\begin{equation}
\ev{u,\SubstrateOperator{\sigma}(y)u}
=
\ev{
\CanonicalLift{\sigma}(y)v,
\SubstrateOperator{\sigma}(y)\CanonicalLift{\sigma}(y)v
}
+
\ev{k,\SubstrateOperator{\sigma}(y)k}.
\label{eq:split}
\end{equation}
Consequently,
\[
\SubstrateEnergy{\sigma}(y,u)
\ge
\CoreForcedBodyResponse{\sigma}(y,v),
\]
with equality if and only if \(k=0\).
\end{proposition}

\begin{proof}
The decomposition \(u=\CanonicalLift{\sigma}(y)v+k\) is Proposition \ref{prop:canonical}(3). Expanding the quadratic form and using Proposition \ref{prop:canonical}(2) gives \eqref{eq:split}. The coercive bound \eqref{eq:coercive} makes the final term nonnegative, and it vanishes only when \(k=0\).
\end{proof}

\section{Canonical Current-Body Image Core}

\begin{definition}[Canonical current-body image core]
\label{def:imagecore}
Fix \(\sigma\in\{\Car,\Cur,\Hom\}\) and a benchmark-faithful current-body-realizing projection-operator core in the sense of Definition \ref{def:core}. The \emph{canonical current-body image core} \(\CurrentBodyImageCore\) for sector \(\sigma\) consists of:
\begin{enumerate}
    \item the same finite-dimensional Euclidean substrate carrier \(\SubstrateCarrier{\sigma}\);
    \item the same surjective displacement projection \(\SubstrateProjection{\sigma}(y)\);
    \item the same coercive positive self-adjoint substrate-response operator \(\SubstrateOperator{\sigma}(y)\);
    \item the canonical current-body image body
    \begin{equation}
    \CanonicalImageBody{\sigma}(y)
    \equiv
    \CanonicalLift{\sigma}(y)\bigl(\PrePreProtoSectionResidualBody{\sigma}(y)\bigr)
    \subset
    \SubstrateCarrier{\sigma};
    \label{eq:imagebody}
    \end{equation}
    \item the exact-shell image
    \begin{equation}
    \CanonicalImageShell{\sigma}
    \equiv
    \left\{
    \bigl(y,\CanonicalLift{\sigma}(y)0\bigr):
    y\in\ProtoSectionShell{\sigma}
    \right\}
    =
    \left\{
    (y,0):
    y\in\ProtoSectionShell{\sigma}
    \right\};
    \label{eq:imageshell}
    \end{equation}
    \item benchmark faithfulness.
\end{enumerate}
From this smaller core define the image-core response on the fixed current displacement body by
\begin{equation}
\CanonicalImageResponse{\sigma}(y,v)
\equiv
\frac{1}{2}
\ev{
\CanonicalLift{\sigma}(y)v,
\SubstrateOperator{\sigma}(y)\CanonicalLift{\sigma}(y)v
}
\qquad
\text{for }
(y,v)\in\PrePreProtoSectionResidualBody{\sigma}.
\label{eq:imageresponse}
\end{equation}
\end{definition}

\begin{remark}
Definition \ref{def:imagecore} is strictly smaller than Definition \ref{def:core}. The full substrate body is no longer kept as an independently primitive compact convex set over the carrier. What is retained benchmark-facingly is only the canonical image of the already-fixed current displacement body under the already-forced minimal lift.
\end{remark}

\begin{proposition}[The canonical image body is exactly the body seen by the current displacements]
\label{prop:reconstructbody}
For every \(y\in\ProtoSectionBody{\sigma}\),
\begin{enumerate}
    \item \(\CanonicalLift{\sigma}(y)\) maps \(\PrePreProtoSectionResidualBody{\sigma}(y)\) bijectively onto \(\CanonicalImageBody{\sigma}(y)\);
    \item the restricted projection
    \[
    \SubstrateProjection{\sigma}(y)\big|_{\CanonicalImageBody{\sigma}(y)}
    :
    \CanonicalImageBody{\sigma}(y)\to\PrePreProtoSectionResidualBody{\sigma}(y)
    \]
    is the inverse of that bijection;
    \item each \(\CanonicalImageBody{\sigma}(y)\) is nonempty, compact, and convex.
\end{enumerate}
\end{proposition}

\begin{proof}
Item (1) is immediate from \eqref{eq:imagebody}. For item (2), if \(u=\CanonicalLift{\sigma}(y)v\in\CanonicalImageBody{\sigma}(y)\), Proposition \ref{prop:canonical}(1) gives
\[
\SubstrateProjection{\sigma}(y)u
=
\SubstrateProjection{\sigma}(y)\CanonicalLift{\sigma}(y)v
=
v.
\]
Conversely, if \(v\in\PrePreProtoSectionResidualBody{\sigma}(y)\), then by definition \(\CanonicalLift{\sigma}(y)v\in\CanonicalImageBody{\sigma}(y)\). Therefore the restricted projection is inverse to the canonical lift on the current displacement body. Item (3) follows because \(\PrePreProtoSectionResidualBody{\sigma}(y)\) is nonempty, compact, and convex and \(\CanonicalLift{\sigma}(y)\) is linear and continuous.
\end{proof}

\begin{proposition}[The exact-shell image is forced by the current-body compatibility relation]
\label{prop:reconstructshell}
For every benchmark-faithful projection-operator core,
\[
\CanonicalImageShell{\sigma}
=
\left\{
(y,u)\in\ProtoSectionBody{\sigma}\times\CanonicalImageBody{\sigma}(y):
\SubstrateProjection{\sigma}(y)u=0
\right\}.
\]
\end{proposition}

\begin{proof}
If \((y,u)\in\CanonicalImageShell{\sigma}\), then by \eqref{eq:imageshell} we have \(u=\CanonicalLift{\sigma}(y)0\), so Proposition \ref{prop:canonical}(1) gives \(\SubstrateProjection{\sigma}(y)u=0\). Conversely, suppose \(u\in\CanonicalImageBody{\sigma}(y)\) and \(\SubstrateProjection{\sigma}(y)u=0\). By Proposition \ref{prop:reconstructbody}(1), \(u=\CanonicalLift{\sigma}(y)v\) for some \(v\in\PrePreProtoSectionResidualBody{\sigma}(y)\). Applying the projection and using Proposition \ref{prop:canonical}(1) gives \(v=0\). Proposition \ref{prop:compat} then implies \(y\in\ProtoSectionShell{\sigma}\), so \((y,u)=(y,\CanonicalLift{\sigma}(y)0)\in\CanonicalImageShell{\sigma}\).
\end{proof}

\begin{proposition}[The current body-response core and downstream chain factor through the image core]
\label{prop:downstream}
For every benchmark-faithful current-body-realizing projection-operator core, the current body-response core and all current benchmark-facing downstream uses of the projection-operator-core frontier already factor through the canonical current-body image core \(\CurrentBodyImageCore\). More precisely:
\begin{enumerate}
    \item
    \[
    \CanonicalImageResponse{\sigma}(y,v)
    =
    \CoreForcedBodyResponse{\sigma}(y,v)
    \]
    for every \((y,v)\in\PrePreProtoSectionResidualBody{\sigma}\);
    \item the pair
    \[
    \left(
    \PrePreProtoSectionResidualBody{\sigma},
    \CanonicalImageResponse{\sigma}
    \right)
    \]
    is the same benchmark-faithful observer-relevant pullback body-response core already forced by the current projection-operator-core theorem;
    \item the downstream pullback body-operator core, pullback-operator core, pullback-quadratic core, section-centered residual-normal-form realization, pre-proto-section package, minimizer-section core, continuity-Hessian compressed core, and benchmark one-metric observer package therefore follow unchanged once \(\CurrentBodyImageCore\) is fixed.
\end{enumerate}
\end{proposition}

\begin{proof}
Item (1) is just the comparison of Definitions \ref{def:canonical} and \ref{def:imagecore}. Item (2) follows because the current displacement body is fixed in Definition \ref{def:bodyresp} and the response on that body is exactly the same forced quadratic form. Item (3) is then exactly the recorded downstream fact of Definition \ref{def:downstream}.
\end{proof}

\begin{proposition}[Off-image substrate-body directions are benchmark neutral]
\label{prop:neutral}
Fix a benchmark-faithful current-body-realizing projection-operator core. For every \(y\in\ProtoSectionBody{\sigma}\) and every \(u\in\SubstrateBody{\sigma}(y)\), if
\[
v\equiv\SubstrateProjection{\sigma}(y)u
\in
\PrePreProtoSectionResidualBody{\sigma}(y),
\qquad
k\equiv u-\CanonicalLift{\sigma}(y)v,
\]
then
\begin{enumerate}
    \item \(\CanonicalLift{\sigma}(y)v\in\CanonicalImageBody{\sigma}(y)\);
    \item \(k\in\SubstrateKernel{\sigma}(y)\);
    \item the canonical image point and the original body point project to the same current displacement \(v\);
    \item the excess contribution of \(u\) beyond the canonical image point is exactly the nonnegative kernel-energy term
    \[
    \frac{1}{2}\ev{k,\SubstrateOperator{\sigma}(y)k};
    \]
    \item varying \(k\) does not alter the current body-response core or any downstream benchmark-facing consequence recorded in Definition \ref{def:downstream}.
\end{enumerate}
\end{proposition}

\begin{proof}
Items (1)--(3) are immediate from the definition of \(\CanonicalImageBody{\sigma}(y)\) and Proposition \ref{prop:canonical}(3). Item (4) is Proposition \ref{prop:split}. For item (5), the body-response core and every listed downstream consequence factor through \(\CanonicalImageResponse{\sigma}\) by Proposition \ref{prop:downstream}; that response depends only on the canonical image representative of each current displacement, not on additional kernel directions inside the larger substrate body.
\end{proof}

\begin{proposition}[The benchmark package remains fixed]
\label{prop:benchmark}
Every canonical current-body image core preserves the same benchmark one-metric package \eqref{eq:fixedoutputs}. In particular, the operative metric remains exactly \(\CompletedMetric\), the ordinary matter route is unchanged, there is no second operative metric, and the low-energy matter law is not enlarged.
\end{proposition}

\begin{proof}
This is part of benchmark faithfulness in Definition \ref{def:imagecore}(6). The present note only removes benchmark-neutral body directions inside the current projection-operator-core frontier.
\end{proof}

\section{Main Theorem}

\begin{theorem}[Projection-operator core collapse to a canonical current-body image core]
\label{thm:main}
Fix the primitive continuation data of Definition \ref{def:fixed}, the recorded proto-section benchmark base and current body-response datum of Definition \ref{def:bodyresp}, the recorded downstream chain of Definition \ref{def:downstream}, and for each sector \(\sigma\in\{\Car,\Cur,\Hom\}\) a benchmark-faithful current-body-realizing projection-operator core in the sense of Definition \ref{def:core}. Then the live benchmark-facing burden inside the current projection-operator-core frontier collapses from the full substrate-body presentation of \(\ProjectionOperatorCore\) to the smaller canonical current-body image core \(\CurrentBodyImageCore\). More precisely:
\begin{enumerate}
    \item the exact-shell/current-body compatibility relation is already forced by the full projection-operator core by Proposition \ref{prop:compat};
    \item the canonical image body and exact-shell image are recovered from the fixed current displacement body by Propositions \ref{prop:reconstructbody} and \ref{prop:reconstructshell};
    \item the current body-response core and all current downstream benchmark-facing consequences factor through \(\CurrentBodyImageCore\) by Proposition \ref{prop:downstream};
    \item every point of the larger substrate body decomposes uniquely into the canonical image representative of the same current displacement together with a kernel direction, and those off-image kernel directions are benchmark neutral at current scope by Proposition \ref{prop:neutral};
    \item the benchmark boundary remains fixed by Proposition \ref{prop:benchmark};
    \item consequently the remaining honest burden below the previous projection-operator-core frontier is no longer the full current-body-realizing substrate body inside \(\ProjectionOperatorCore\) as such. What remains is to derive or justify the smaller canonical current-body image core \(\CurrentBodyImageCore\) from still more primitive explicit substrate structure, or else to prove a still smaller collapse, sharper necessity result, sharper uniqueness result, or sharper compatibility theorem strictly inside that smaller image core.
\end{enumerate}
\end{theorem}

\begin{proof}
Item (1) is Proposition \ref{prop:compat}. Item (2) is the combined content of Propositions \ref{prop:reconstructbody} and \ref{prop:reconstructshell}. Item (3) is Proposition \ref{prop:downstream}. Item (4) is Proposition \ref{prop:neutral}. Item (5) is Proposition \ref{prop:benchmark}. Item (6) is the routing consequence of items (1)--(5) together with the active safeguard against counting another same-output deeper relay as new front-line progress when the live benchmark-relevant datum is unchanged.
\end{proof}

\begin{corollary}[Primary post-collapse burden]
\label{cor:next}
After Theorem \ref{thm:main}, the next honest benchmark-facing manuscript below the present frontier must do at least one of the following:
\begin{enumerate}
    \item derive or force the canonical current-body image core \(\CurrentBodyImageCore\) itself from still more primitive explicit substrate structure;
    \item derive or force a nontrivial constituent or relation inside \(\CurrentBodyImageCore\), such as the canonical image body, the restricted projection/canonical lift relation, the exact-shell image, or the response operator restricted to the canonical image;
    \item prove a sharper necessity, uniqueness, or compatibility theorem for some nontrivial constituent or relation inside that smaller image core;
    \item do one of the previous three while preserving the same benchmark one-metric package, the same ordinary matter route, and the same claim boundary.
\end{enumerate}
Another manuscript that merely reinstates the full substrate body as if off-image kernel directions were still independently live, merely replays the full projection-operator core, merely returns to the full displacement-response substrate law, or re-expands back to the larger pullback-body-operator, pullback-operator, pullback-quadratic, hidden-response, residual-normal-form, or pre-proto-section presentations is not the honest default next move.
\end{corollary}

\begin{proof}
Theorem \ref{thm:main} shows that the larger substrate-body presentation is no longer the minimal benchmark-facing datum at this stage. Therefore a subsequent manuscript must improve on the smaller image core rather than merely restating the larger projection-operator core.
\end{proof}

\section{Conclusion}

The positive result of this note is a genuine smaller theorem strictly inside the current projection-operator-core frontier. The previous active-primary theorem already showed that the full section-centered displacement-response substrate law is not the minimal benchmark-facing datum because the separately primitive minimal-lift family collapses to a smaller current-body-realizing projection-operator core. The present theorem then takes the honest next internal route available there: it shows that even that smaller core still contains benchmark-neutral substrate-body freedom.

What survives as the live burden is the canonical current-body image core \(\CurrentBodyImageCore\). Relative to the already fixed current displacement body, the only body points that matter at current benchmark-facing scope are the canonical minimal-lift representatives of those current displacements. Any additional substrate-body direction is a kernel direction that contributes only nonnegative excess response energy and leaves the forced body-response core and every recorded downstream benchmark-facing consequence unchanged.

The scope remains honest. The one operative metric is still exactly \(\CompletedMetric\), the ordinary matter route is unchanged, there is no second operative metric, and the low-energy matter law is not enlarged. The current result does not complete a first-principles derivation of GR from substrate microphysics, but it does narrow the live derivational burden one level further on the active line. The next open burden is therefore to derive or justify \(\CurrentBodyImageCore\) itself from still more primitive explicit substrate structure, or else to prove a still sharper theorem strictly inside that smaller image core.

\end{document}
